% Pursuing Stacks (cotes 134-2 to 134-8): archive cote 134-3, site page 83 (PDF page 84), typescript page 83.
% typescript. Draft compared: ps2.tex:4326-4391.
% AI-transcribed by Claude Opus 5.5 (claude-opus-5-5, adaptive thinking, effort high), Message Batches API, batch ps1 (msgbatch_01Rv5hCaNTtUt4JHx1yqoHPL), 2026-11-12; image review recorded in records/CHECK_DECISIONS_134_3_4.tsv.
\scanpage{134-3/83}
\mspage{83}

final object)  over $e_A = e_{A^{\hand{\wedge}}}$ .

We express the conditions (i) to (v)  by saying that A is a \uline{local}
\uline{test} \uline{category} ,  or that A\hand{\ensuremath{^{\wedge}}}\hand{/} is a \uline{locally modelizing topos},  or (if we want
to recall that this topos is of the  type A\hand{\ensuremath{^{\wedge}}}) an \uline{elementary local modeliz\xout{ze}}
\uline{zer}.  If A is moreover aspheric  (i.e. $e_{A^{\hand{\wedge}}}$ aspheric),  or what amounts to
the same,  if A is moreover a weak test category,  we\deleted{\xout{'ll}} say that A is a
\uline{test} \uline{category} ,  or that A\hand{\ensuremath{^{\wedge}}} is a \uline{modelizing} \uline{topos}, or (to recall it is an
A\hand{\ensuremath{^{\wedge}}} and not just any topos) an \uline{elementary} \uline{modelizer}.  These are intrinsic
properties on the topos A\hand{\ensuremath{^{\wedge}}},  the first one of a local nature,  the second
not (as asphericity of A\hand{\ensuremath{^{\wedge}}} is a global notion).

Here the question arises whether the condition for A to be a weak \xout{tes}
test category can be likewise expressed intrinsically as a property of the
topos A\hand{\ensuremath{^{\wedge}}} -  which we then would call a \uline{weakly modelizing topos} ,  or a
\uline{weak elementary modelizer}.  This does'nt look so clear,  as all conditions
stated in a)  make use of at least one among the two functors $i_A$ , $j_A$ ,
which do not seem to make much sense \inserted{in the more general case}, except \xout{when} possibly when using a
specified small generating subcategory A of the given topos   -  and
possibly checking that the condition\xout{s} obtained (if something reasonable
comes out,  as I do expect) does not depend on the choice of the genera-%
ting site.  This should be part of a systematic reflection on modelizing
topoi,  \xout{with} to make sure \inserted[typed]{for instance} they are \xout{\unsure{inde}} modelizing indeed with respect to
weak equivalences  -  \over{and}{but} I'll not \xout{\unsure{dwell on}} enter into such reflection
for the time being.

c)   \uline{Strict test categories}.   The following conditions \over[typed]{o}{o}\over[typed]{n}{n} the small cate-%
gory A are equivalent:

(i)   A is a weak test category \inserted[typed]{(cf.  a) (iii) above)},  and thus induces a localization
functor
\[ \deleted{\illegible{1}}\ A^{\hand{\wedge}} \hand{\xrightarrow{\deleted{\illegible{1}}}} (\mathrm{Hot}) = W^{-1}(\mathrm{Cat}) \quad , \]
and this functor moreover \uline{commutes with binary products};  or equi-%
valently,  the canonical functor $A^{\hand{\wedge}} \hand{\to} W_A^{-1}A^{\hand{\wedge}}$\hand{/} commutes with binary
products.

(ii)  A is a test category (a local test category would be enough
ev\unsure{en}),  and moreover the topos A\hand{\ensuremath{^{\wedge}}} is totally aspheric \inserted{[Section 35]}, namely
(apart \inserted{from} having a generating family of aspheric generators,  which
is clear anyhow)  the product of any two aspheric elements in A\hand{\ensuremath{^{\wedge}}}
is aspheric,  i.e. any aspheric element of A\hand{\ensuremath{^{\wedge}}}\hand{/} is aspheric over the
final object.
